% ============================================================
%  APUNTES — Bloque 2: Templates Modernos (Angular 17+)
%  Compilar con: LuaLaTeX
%  En Overleaf: Menu -> Compiler -> LuaLaTeX
% ============================================================
\documentclass[12pt, a4paper]{article}

% --- Fuentes y codificación (LuaLaTeX) ----------------------
\usepackage{fontspec}
\setmainfont{Latin Modern Roman}
\setmonofont[Scale=0.88]{Courier New}

% --- Emojis -------------------------------------------------
\usepackage{emoji}
\setemojifont{Segoe UI Emoji}

% --- Idioma -------------------------------------------------
\usepackage[spanish]{babel}

% --- Márgenes -----------------------------------------------
\usepackage[top=2.5cm, bottom=2.5cm, left=2.8cm, right=2.8cm]{geometry}

% --- Colores ------------------------------------------------
\usepackage[dvipsnames, table]{xcolor}
\definecolor{azulSignal}{RGB}{37, 99, 235}
\definecolor{verdeOk}{RGB}{22, 163, 74}
\definecolor{rojoError}{RGB}{220, 38, 38}
\definecolor{grisClaro}{RGB}{243, 244, 246}
\definecolor{grisCodigo}{RGB}{31, 41, 55}
\definecolor{amarillo}{RGB}{251, 191, 36}
\definecolor{morado}{RGB}{124, 58, 237}

% --- Código fuente ------------------------------------------
\usepackage{listings}
\lstdefinestyle{ts}{
  language        = [Sharp]C,
  basicstyle      = \ttfamily\small,
  keywordstyle    = \color{azulSignal}\bfseries,
  commentstyle    = \color{verdeOk}\itshape,
  stringstyle     = \color{rojoError},
  backgroundcolor = \color{grisClaro},
  frame           = single,
  framerule       = 0.5pt,
  rulecolor       = \color{gray!40},
  numbers         = left,
  numberstyle     = \tiny\color{gray},
  numbersep       = 8pt,
  tabsize         = 2,
  breaklines      = true,
  showstringspaces = false,
  xleftmargin     = 1.5em,
  morekeywords    = {if, else, for, switch, case, default, let,
                     defer, placeholder, loading, error, empty,
                     track, on, when, signal, computed, input,
                     output, model, readonly, const},
}
\lstset{style=ts}

% --- Cajas decorativas --------------------------------------
\usepackage[most]{tcolorbox}
\tcbset{
  analogia/.style={
    colback=blue!5, colframe=azulSignal,
    fonttitle=\bfseries\large, arc=6pt,
    left=8pt, right=8pt, top=6pt, bottom=6pt,
  },
  importante/.style={
    colback=yellow!10, colframe=amarillo!80!orange,
    fonttitle=\bfseries\large, arc=6pt,
    left=8pt, right=8pt,
  },
  regla/.style={
    colback=green!5, colframe=verdeOk,
    fonttitle=\bfseries\large, arc=6pt,
    left=8pt, right=8pt,
  },
  peligro/.style={
    colback=red!5, colframe=rojoError,
    fonttitle=\bfseries\large, arc=6pt,
    left=8pt, right=8pt,
  },
}

% --- Resto de paquetes --------------------------------------
\usepackage{hyperref}
\hypersetup{colorlinks=true, linkcolor=azulSignal, urlcolor=morado}
\usepackage{booktabs}
\usepackage{array}
\usepackage{fancyhdr}
\pagestyle{fancy}
\fancyhf{}
\rhead{\textcolor{gray}{\small \emoji{artist-palette} Templates --- Angular 17+}}
\lhead{\textcolor{gray}{\small Survey Maker}}
\cfoot{\textcolor{gray}{\thepage}}
\setlength{\parskip}{6pt}
\setlength{\parindent}{0pt}

% ============================================================
\begin{document}

% --- Portada ------------------------------------------------
\begin{titlepage}
  \centering
  \vspace*{2.5cm}
  {\fontsize{48}{48}\selectfont \emoji{artist-palette}}\\[0.4cm]
  {\Huge\bfseries\color{azulSignal} Apuntes de Angular}\\[0.4cm]
  {\Large\color{grisCodigo} Bloque 2 --- Templates Modernos (Angular 17+)}\\[0.8cm]
  \textcolor{gray!60}{\rule{10cm}{0.8pt}}\\[0.8cm]
  {\large Proyecto: \textbf{Survey Maker}}\\[0.2cm]
  {\normalsize\texttt{\textcolor{morado}{github.com/AldoZM/Survey-Maker}}}\\[0.8cm]
  {\normalsize \emoji{student} Aldo Zetina --- UNAM Ciencias de la Computación}\\[0.2cm]
  {\normalsize \emoji{calendar} Mayo 2026}\\[2cm]
  \begin{tcolorbox}[importante, title={\emoji{wrench} Compilar en Overleaf}]
    \textbf{Menu} \textrightarrow\ \textbf{Compiler} \textrightarrow\ \textbf{LuaLaTeX}
  \end{tcolorbox}
\end{titlepage}

\tableofcontents
\newpage

% ============================================================
\section{\emoji{brain} ¿Qué es un Template?}

El template es el HTML que define \textbf{lo que el usuario ve}.
Pero no es HTML normal --- Angular lo procesa y lo hace dinámico.

\begin{tcolorbox}[analogia, title={\emoji{cook} La Receta con Instrucciones}]
Piensa en el template como una \textbf{receta de cocina con instrucciones condicionales}:

\textit{``Si hay huevos, haz omelette. Si no hay, haz cereal.
Para cada ingrediente en la lista, agrégalo al plato.''}

Eso es exactamente lo que hacen \texttt{@if} y \texttt{@for} ---
le das instrucciones a Angular y él decide qué renderizar
según el estado actual de los Signals.
\end{tcolorbox}

% ============================================================
\section{\emoji{speech-balloon} Interpolación \texttt{\{\{ \}\}}}

La forma más simple de mostrar datos en pantalla.
Evalúa una expresión TypeScript y la convierte a texto visible.

\begin{lstlisting}
<!-- Muestra el titulo del schema -->
<mat-card-title>{{ schema().title }}</mat-card-title>
<!--                ^
                    schema() es Signal — () obligatorios para leer -->

<!-- Operaciones dentro de {{ }} -->
<span>{{ state.progress() }}% completado</span>
<!-- Si progress() = 60 renderiza: "60% completado" -->

<!-- Operador ?? para valores nulos -->
<button>{{ schema().submitLabel ?? 'Enviar' }}</button>
<!--                              ^
                    submitLabel nulo/undefined -> muestra 'Enviar' -->
\end{lstlisting}

% ============================================================
\section{\emoji{link} Property Binding \texttt{[ ]}}

Conecta una expresión TypeScript a una \textbf{propiedad} del elemento HTML
o de un componente hijo. La diferencia con el texto plano es crítica:

\begin{lstlisting}
<!-- SIN corchetes — pasa el TEXTO literal "value()" -->
<input value="value()">
<!-- El input muestra la cadena "value()" — INCORRECTO -->

<!-- CON corchetes — evalua la EXPRESION y pasa el resultado -->
<input [value]="value()">
<!-- El input muestra lo que retorne value() — CORRECTO -->

<!-- Mas ejemplos del proyecto -->
<mat-progress-bar
  [value]="state.progress()"   <!-- evalua la expresion Signal -->
  mode="determinate"           <!-- string fijo, sin corchetes -->
/>

<mat-checkbox
  [checked]="isChecked(option.value)"  <!-- true/false dinamico -->
  [disabled]="option.disabled ?? false"
/>
\end{lstlisting}

% ============================================================
\section{\emoji{ear} Event Binding \texttt{( )}}

Escucha \textbf{eventos} del DOM o de componentes hijos y ejecuta
código cuando ocurren.

\begin{lstlisting}
<!-- Eventos nativos del browser -->
<input (input)="onInput($event)">
<!--    ^       ^
        evento  funcion que corre cuando ocurre
        $event = objeto Event nativo del browser -->

<button (click)="onSubmit()">Enviar</button>

<!-- Outputs de componentes hijos (custom events) -->
<app-select-field
  (valueChange)="state.setAnswer(field.id, $event)"
  <!--            ^
                  output del hijo emitido con valueChange.emit(x)
                  $event = ese valor x -->
/>

<!-- Modificador de tecla — solo dispara con Enter -->
<input (keyup.enter)="buscar()">
\end{lstlisting}

\begin{tcolorbox}[regla, title={\emoji{pushpin} La Regla de los Corchetes}]
  \textbf{Datos bajan con \texttt{[ ]}} --- del padre al hijo.\\
  \textbf{Eventos suben con \texttt{( )}} --- del hijo al padre.\\[0.3cm]
  Juntos: \texttt{[( )]} = two-way binding (los dos en uno).
\end{tcolorbox}

% ============================================================
\section{\emoji{performing-arts} \texttt{@if} / \texttt{@else} --- Renderizado Condicional}

\subsection{La Historia: De \texttt{*ngIf} a \texttt{@if}}

\begin{lstlisting}
<!-- VIEJO — *ngIf, el else era confuso y verboso -->
<div *ngIf="condicion; else miElse">Si</div>
<ng-template #miElse><div>No</div></ng-template>

<!-- NUEVO Angular 17+ — limpio como TypeScript normal -->
@if (condicion) {
  <div>Si</div>
} @else {
  <div>No</div>
}
\end{lstlisting}

\subsection{En tu proyecto --- \texttt{survey-page.component.ts}}

\begin{lstlisting}
<!-- Estado 1: Descargando el JSON de la encuesta -->
@if (loader.survey.isLoading()) {
  <!-- Solo existe en el DOM mientras isLoading() = true -->
  <!-- Cuando termina, Angular ELIMINA este nodo del DOM -->
  <!-- No es display:none — literalmente no existe en memoria -->
  <mat-progress-bar mode="indeterminate" />
}

<!-- Estado 2: Algo fallo (404, red caida, JSON invalido) -->
@if (loader.survey.error(); as err) {
  <!--                      ^
                            "as err" captura el valor del Signal
                            en variable local para usarla dentro
                            Sin "as", tendrias que escribir
                            loader.survey.error() dos veces -->
  <div class="error-container">
    <p>Error al cargar la encuesta.</p>
    <pre>{{ err }}</pre>
  </div>
}

<!-- Estado 3: Llego el SurveySchema validado -->
@if (loader.survey.value(); as schema) {
  <!--                      ^
                            "as schema" evita escribir
                            [schema]="loader.survey.value()!"
                            con el ! de non-null assertion -->
  <app-survey [schema]="schema" />
}
\end{lstlisting}

\begin{tcolorbox}[peligro, title={\emoji{warning} \texttt{@if} NO es \texttt{display:none}}]
\begin{center}
\renewcommand{\arraystretch}{1.4}
\begin{tabular}{p{5.5cm} p{6cm}}
\textbf{\texttt{display:none} (CSS)} & \textbf{\texttt{@if} (Angular)} \\
\hline
El elemento \textbf{existe} en el DOM & El elemento \textbf{no existe} \\
Angular ejecuta su código & Angular no ejecuta nada \\
Consume memoria y CPU & Cero costo cuando es \texttt{false} \\
No dispara \texttt{ngOnDestroy()} & Dispara \texttt{ngOnDestroy()} \\
\end{tabular}
\end{center}
\end{tcolorbox}

% ============================================================
\section{\emoji{repeat-button} \texttt{@for} --- Iteración Reactiva}

\subsection{La Historia: De \texttt{*ngFor} a \texttt{@for}}

\begin{lstlisting}
<!-- VIEJO — trackBy era opcional, muchos lo omitian -->
<div *ngFor="let s of sections; trackBy: trackByFn">
  {{ s.title }}
</div>

<!-- NUEVO — track es OBLIGATORIO (Angular te fuerza a hacerlo bien) -->
@for (section of schema().sections; track section.id) {
  <div>{{ section.title }}</div>
}
\end{lstlisting}

\subsection{Por Qué \texttt{track} es Crítico}

\begin{tcolorbox}[analogia, title={\emoji{id-button} La Huella Digital de Cada Elemento}]
Imagina esta lista en pantalla: \texttt{["Manzana", "Banana", "Cereza"]}\\
El usuario filtra y queda: \texttt{["Banana", "Cereza"]}\\[0.3cm]

\textbf{Sin \texttt{track}:} Angular no sabe cuál desapareció.
Destruye los 3 y recrea 2. Animaciones rotas, focus perdido, inputs limpios.\\[0.3cm]

\textbf{Con \texttt{track section.id}:} Angular sabe que ``Manzana'' desapareció.
Solo destruye ese. Los otros 2 siguen intactos en el DOM.
\end{tcolorbox}

\subsection{En tu proyecto --- \texttt{survey.component.ts}}

\begin{lstlisting}
@for (section of schema().sections; track section.id) {
  <section class="survey-section">
    <h3>{{ section.title }}</h3>

    <!-- @for anidado — campos dentro de cada seccion -->
    @for (field of section.fields; track field.id) {
      <app-field-renderer
        [field]="field"
        [answers]="answersAsValues()"
        [value]="state.answers()[field.id]?.value"
        [errorMessage]="state.getError(field.id)"
        (valueChange)="state.setAnswer(field.id, $event)"
      />
    }
  </section>
}
\end{lstlisting}

\subsection{Variables Automáticas de \texttt{@for}}

\begin{lstlisting}
@for (item of lista; track item.id;
      let i = $index, first = $first, last = $last) {
  <!--      ^              ^               ^
            posicion 0,1,2 true si primero  true si ultimo
            Otras: $even/$odd (par/impar), $count (total) -->

  <div [class.ultimo]="last">
    {{ i + 1 }}. {{ item.nombre }}
    @if (first) { <span>Primero!</span> }
  </div>
}
\end{lstlisting}

\subsection{\texttt{@empty} --- Cuando la Lista está Vacía}

\begin{lstlisting}
@for (encuesta of encuestas; track encuesta.id) {
  <tarjeta [data]="encuesta" />
} @empty {
  <!-- Solo se renderiza cuando encuestas.length === 0 -->
  <p>No hay encuestas disponibles todavia.</p>
}
\end{lstlisting}

% ============================================================
\section{\emoji{left-right-arrow} \texttt{@switch} --- Despacho por Tipo}

Selecciona qué renderizar según el valor de una expresión.
Solo existe en el DOM el \texttt{@case} que coincide.

\begin{lstlisting}
@let f = field();
<!-- @let crea variable local — evita llamar field() 7 veces -->

@switch (f.type) {
  @case ('text') {
    <app-text-field [field]="f" [value]="stringValue()"
      (valueChange)="valueChange.emit($event)" />
  }
  @case ('email') {
    <!-- email reutiliza TextFieldComponent — cambia type="email" interno -->
    <app-text-field [field]="f" [value]="stringValue()"
      (valueChange)="valueChange.emit($event)" />
  }
  @case ('number') {
    <!-- numberValue() garantiza number | null, no string -->
    <app-number-field [field]="f" [value]="numberValue()"
      (valueChange)="valueChange.emit($event)" />
  }
  @case ('select') {
    <app-select-field [field]="f" [value]="stringValue()"
      (valueChange)="valueChange.emit($event)" />
  }
  @case ('radio') {
    <app-radio-field [field]="f" [value]="stringValue()"
      (valueChange)="valueChange.emit($event)" />
  }
  @case ('checkbox') {
    <!-- arrayValue() garantiza string[] para seleccion multiple -->
    <app-checkbox-field [field]="f" [value]="arrayValue()"
      (valueChange)="valueChange.emit($event)" />
  }
}
\end{lstlisting}

% ============================================================
\section{\emoji{pencil} \texttt{@let} --- Variables Locales en Template}

Nueva en Angular 18. Evita repetir expresiones largas y costosas.

\begin{lstlisting}
<!-- Sin @let — expresion larga repetida tres veces -->
@if (state.answers()[field.id]?.value) {
  <span>{{ state.answers()[field.id]?.value }}</span>
  <div [class.filled]="state.answers()[field.id]?.value !== null">
    ...
  </div>
}

<!-- Con @let — limpio y eficiente, una sola evaluacion -->
@let ans = state.answers()[field.id];
@if (ans?.value) {
  <span>{{ ans.value }}</span>
  <div [class.filled]="ans.value !== null">...</div>
}
\end{lstlisting}

% ============================================================
\section{\emoji{hourglass-not-done} \texttt{@defer} --- Lazy Loading en Template}

La novedad más potente de Angular 17. Carga componentes pesados
\textbf{solo cuando los necesitas}, directamente desde el template
sin tocar las rutas.

\begin{lstlisting}
@defer (on interaction) {
  <!-- Se descarga SOLO cuando el usuario hace click/hover -->
  <editor-de-texto-gigante [contenido]="post" />
} @placeholder {
  <!-- Se muestra ANTES de interactuar -->
  <div class="placeholder">Haz clic para editar...</div>
} @loading (minimum 300ms) {
  <!-- Se muestra MIENTRAS se descarga el chunk JS -->
  <mat-spinner diameter="32" />
} @error {
  <!-- Se muestra si la descarga falla (sin red, chunk roto) -->
  <p>Error al cargar el editor. Recarga la pagina.</p>
}
\end{lstlisting}

\subsection{Triggers de \texttt{@defer}}

\begin{center}
\renewcommand{\arraystretch}{1.5}
\begin{tabular}{>{\ttfamily}p{4.5cm} p{8cm}}
\toprule
\textbf{Trigger} & \textbf{Cuándo carga} \\
\midrule
on idle & Cuando el browser no tiene nada que hacer \\
on viewport & Cuando el elemento entra en la pantalla visible \\
on interaction & Cuando el usuario hace hover o click \\
on hover & Solo al pasar el mouse encima \\
on timer(3s) & Después de 3 segundos automáticamente \\
when condicion() & Cuando una expresión Signal es \texttt{true} \\
\bottomrule
\end{tabular}
\end{center}

% ============================================================
\section{\emoji{banana} Two-Way Binding \texttt{[( )]}}

Llamado \textit{``banana in a box''} porque \texttt{[()]} parece un
plátano dentro de una caja. Combina \texttt{[]} y \texttt{()} en uno.

\begin{lstlisting}
<!-- Sin two-way binding — verboso, dos atributos -->
<input [value]="nombre" (input)="nombre = $event.target.value" />

<!-- Con two-way binding — conciso, un solo atributo -->
<input [(ngModel)]="nombre" />

<!-- Con el nuevo model() Signal (Angular 17+) -->
export class MiComponente {
  readonly nombre = model<string>('');
  // Al hacer [(nombre)]="miVar" en el padre:
  //   [nombre]       → pasa miVar al componente
  //   (nombreChange) → cuando nombre.set() corre, actualiza miVar
}
\end{lstlisting}

% ============================================================
\section{\emoji{abacus} Pipes --- Transformadores en el Template}

Los pipes transforman valores \textbf{sin modificar el dato original}.
Son como filtros de Instagram: no cambian la foto, solo cómo se ve.

\begin{lstlisting}
{{ fecha | date:'dd/MM/yyyy' }}
<!-- 2026-05-13 -> "13/05/2026" -->

{{ precio | currency:'MXN' }}
<!-- 1500 -> "MX$1,500.00" -->

{{ texto | uppercase }}
<!-- "hola mundo" -> "HOLA MUNDO" -->

{{ objeto | json }}
<!-- util para debugging: { id: 1 } -> '{"id":1}' -->

<!-- Pipes encadenados de izquierda a derecha -->
{{ nombre | uppercase | slice:0:10 }}
<!-- "Aldo Zetina Meza" -> "ALDO ZETIN" -->
\end{lstlisting}

% ============================================================
\section{\emoji{clipboard} Resumen del Bloque 2}

\begin{center}
\renewcommand{\arraystretch}{1.6}
\begin{tabular}{>{\ttfamily\bfseries}p{4.5cm} p{8cm}}
\toprule
\textbf{Sintaxis} & \textbf{Para qué} \\
\midrule
\{\{ expr \}\} & Mostrar texto en pantalla (interpolación) \\
{[}prop{]} & Pasar dato a elemento o componente hijo \\
(evento) & Escuchar evento de elemento o componente hijo \\
{[}(modelo){]} & Two-way binding (los dos anteriores juntos) \\
\midrule
@if (c) \{ \} & Renderizar si condición es \texttt{true} \\
@else if / @else \{ \} & Alternativas al \texttt{@if} \\
@for (x of l; track x.id) & Iterar lista de forma eficiente \\
@empty \{ \} & Cuando la lista está vacía \\
@switch / @case / @default & Seleccionar según valor \\
@let var = expr; & Variable local reutilizable en template \\
@defer (trigger) \{ \} & Carga diferida de componentes pesados \\
\bottomrule
\end{tabular}
\end{center}

\vfill
\begin{center}
  \textcolor{gray}{\small \emoji{artist-palette} Apuntes de sesión --- Mayo 2026}\\[0.2cm]
  \textcolor{gray}{\small Survey Maker --- Angular 19 Signal-First}\\[0.2cm]
  \textcolor{gray}{\small \texttt{github.com/AldoZM/Survey-Maker}}
\end{center}

\end{document}
